
\documentclass{article}
\usepackage{colortbl}
\usepackage{makecell}
\usepackage{multirow}
\usepackage{supertabular}

\begin{document}

\newcounter{utterance}

\centering \large Interaction Transcript for game `bbh', experiment `salient\_translation\_error\_detection', episode 55 with reflexion{-}Qwen3.5{-}27B.
\vspace{24pt}

{ \footnotesize  \setcounter{utterance}{1}
\setlength{\tabcolsep}{0pt}
\begin{supertabular}{c@{$\;$}|p{.15\linewidth}@{}p{.15\linewidth}p{.15\linewidth}p{.15\linewidth}p{.15\linewidth}p{.15\linewidth}}
    \# & \multicolumn{2}{c}{Player} && \multicolumn{2}{c}{Game Master} \\
    \hline

    \theutterance \stepcounter{utterance}  
    & & & \multicolumn{4}{p{0.6\linewidth}}{
        \cellcolor[rgb]{0.9,0.9,0.9}{
            \makecell[{{p{\linewidth}}}]{
                \texttt{\tiny{[P1$\langle$GM]}}
                \texttt{You have attempted this same task before and did not succeed. These are the plans you wrote after those attempts.} \\
\texttt{Trial \#1: New plan: In the previous attempt, I correctly identified the error as "Negation or Antonyms" but failed to output the answer in the strict format required by the few{-}shot examples. The examples show that the final answer must be a single line containing only the option letter in parentheses (e.g., "A: (D)"), without any preceding explanation or reasoning text. My mistake was providing a detailed analysis before the final answer. Next time, I will analyze the source and translation to identify the error type internally, but I will output **only** the final line in the format "A: (Option Letter)" to match the pattern established in the few{-}shot examples.} \\
\\ 
\texttt{Trial \#2: New plan: In the previous attempt, I incorrectly classified the error as "Negation or Antonyms" (Option C) because I focused on the change from "unvollständige" (incomplete) to "complete," interpreting it as an antonym swap. However, the source sentence explicitly lists specific districts (Bantorf, Barrigsen, etc.), and the translation includes all of them. The error is not a negation or antonym change, but rather a change in the modifier describing the list itself: "unvollständige" (incomplete) was translated as "complete." This falls under "Modifiers or Adjectives" (Option A), as the adjective describing the noun "Liste" was altered to its opposite meaning, but the specific category for changing adjectives/modifiers is distinct from the "Negation or Antonyms" category which typically involves adding/removing "not" or swapping verbs/nouns to antonyms. More critically, looking at the previous examples, "Negation or Antonyms" was used for "is not" vs "is" (Example 3). Here, the adjective "incomplete" became "complete." While this is an antonym, the category "Modifiers or Adjectives" specifically covers "modifiers and adjectives pertaining to a noun are changed." The change from "incomplete" to "complete" is a direct change of the adjective. I will prioritize checking if the error is a direct adjective swap (Modifiers or Adjectives) before defaulting to Negation/Antonyms, especially when the core noun phrase structure remains but the descriptive adjective is flipped. I will select Option (A) for this case.} \\
\\ 
\\ 
\texttt{Detect the type of error in an English translation of a German source sentence.} \\
\\ 
\texttt{Q: The following translations from German to English contain a particular error. That error will be one of the following types: Named Entities: An entity (names, places, locations, etc.) is changed to a different entity. Numerical Values: Numerical values (ordinals or cardinals), dates, and/or units are changed. Modifiers or Adjectives: The modifiers and adjectives pertaining to a noun are changed. Negation or Antonyms: Introduce or remove a negation or change comparatives to their antonyms. Facts: Trivial factual errors not pertaining to the above classes are introduced in the translations. Dropped Content: A significant clause in the translation is removed. Please identify that error.  Source: In der Liste der Baudenkmale in Lenzen (Elbe) sind alle Baudenkmale der brandenburgischen Stadt Lenzen (Elbe) und ihrer Ortsteile aufgelistet.} \\
\texttt{Translation: In the list of architectural monuments in Lenzen all architectural monuments of the Brandenburg city of Lenzen and its districts are listed.} \\
\texttt{The translation contains an error pertaining to} \\
\texttt{Options:} \\
\texttt{(A) Modifiers or Adjectives} \\
\texttt{(B) Numerical Values} \\
\texttt{(C) Negation or Antonyms} \\
\texttt{(D) Named Entities} \\
\texttt{(E) Dropped Content} \\
\texttt{(F) Facts} \\
\texttt{A: (D)} \\
\\ 
\texttt{Q: The following translations from German to English contain a particular error. That error will be one of the following types: Named Entities: An entity (names, places, locations, etc.) is changed to a different entity. Numerical Values: Numerical values (ordinals or cardinals), dates, and/or units are changed. Modifiers or Adjectives: The modifiers and adjectives pertaining to a noun are changed. Negation or Antonyms: Introduce or remove a negation or change comparatives to their antonyms. Facts: Trivial factual errors not pertaining to the above classes are introduced in the translations. Dropped Content: A significant clause in the translation is removed. Please identify that error.  Source: Auf dieser Seite sind die Baudenkmäler der oberbayerischen Großen Kreisstadt Landsberg am Lech zusammengestellt.} \\
\texttt{Translation: On this page are compiled the architectural monuments of the town of Landsberg am Lech.} \\
\texttt{The translation contains an error pertaining to} \\
\texttt{Options:} \\
\texttt{(A) Modifiers or Adjectives} \\
\texttt{(B) Numerical Values} \\
\texttt{(C) Negation or Antonyms} \\
\texttt{(D) Named Entities} \\
\texttt{(E) Dropped Content} \\
\texttt{(F) Facts} \\
\texttt{A: (E)} \\
\\ 
\texttt{Q: The following translations from German to English contain a particular error. That error will be one of the following types: Named Entities: An entity (names, places, locations, etc.) is changed to a different entity. Numerical Values: Numerical values (ordinals or cardinals), dates, and/or units are changed. Modifiers or Adjectives: The modifiers and adjectives pertaining to a noun are changed. Negation or Antonyms: Introduce or remove a negation or change comparatives to their antonyms. Facts: Trivial factual errors not pertaining to the above classes are introduced in the translations. Dropped Content: A significant clause in the translation is removed. Please identify that error.  Source: Łeba ist eine Kleinstadt und ein Badeort im Powiat Lęborski der polnischen Woiwodschaft Pommern.} \\
\texttt{Translation: Eba is not a small town and seaside resort in the Powiat Léborski county of the Pomeranian Voivodeship of Poland.} \\
\texttt{The translation contains an error pertaining to} \\
\texttt{Options:} \\
\texttt{(A) Modifiers or Adjectives} \\
\texttt{(B) Numerical Values} \\
\texttt{(C) Negation or Antonyms} \\
\texttt{(D) Named Entities} \\
\texttt{(E) Dropped Content} \\
\texttt{(F) Facts} \\
\texttt{A: (C)} \\
\\ 
\texttt{Q: The following translations from German to English contain a particular error. That error will be one of the following types: Named Entities: An entity (names, places, locations, etc.) is changed to a different entity. Numerical Values: Numerical values (ordinals or cardinals), dates, and/or units are changed. Modifiers or Adjectives: The modifiers and adjectives pertaining to a noun are changed. Negation or Antonyms: Introduce or remove a negation or change comparatives to their antonyms. Facts: Trivial factual errors not pertaining to the above classes are introduced in the translations. Dropped Content: A significant clause in the translation is removed. Please identify that error.  Source: Die unvollständige Liste der Baudenkmale in Barsinghausen enthält Baudenkmale der Barsinghausener Kernstadt sowie der Ortsteile Bantorf, Barrigsen, Eckerde, Egestorf, Göxe, Großgoltern, Groß Munzel, Hohenbostel, Holtensen, Kirchdorf, Landringhausen, Langreder, Nordgoltern, Ostermunzel, Stemmen, Wichtringhausen und Winninghausen.} \\
\texttt{Translation: The complete list of architectural monuments in Barsinghausen contains architectural monuments of the Barsinghausen core town as well as the districts bantorf, Barrigsen, Eckerde, Egestorf, Göxe, Großgoltern, Groß Munzel, Hohenbostel, Holtensen, Kirchdorf, Landringhausen, Langreder, Nordgoltern, Ostermunzel, Stemmen, Wichtringhausen and Winninghausen.} \\
\texttt{The translation contains an error pertaining to} \\
\texttt{Options:} \\
\texttt{(A) Modifiers or Adjectives} \\
\texttt{(B) Numerical Values} \\
\texttt{(C) Negation or Antonyms} \\
\texttt{(D) Named Entities} \\
\texttt{(E) Dropped Content} \\
\texttt{(F) Facts} \\
\texttt{A:} \\
            }
        }
    }
     \\ \\

    \theutterance \stepcounter{utterance}  
    & \multicolumn{4}{p{0.6\linewidth}}{
        \cellcolor[rgb]{0.561,0.737,0.561}{
            \makecell[{{p{\linewidth}}}]{
                \texttt{\tiny{[P1$\rangle$GM]}}
                \texttt{A: (A)} \\
            }
        }
    }
    & & \\ \\

    \theutterance \stepcounter{utterance}  
    & & & \multicolumn{2}{p{0.3\linewidth}}{
        \cellcolor[rgb]{0.9,0.9,0.9}{
            \makecell[{{p{\linewidth}}}]{
                \texttt{\tiny{[GM$|$GM]}}
                \texttt{Target: (C)} \\
            }
        }
    }
    & & \\ \\

    \theutterance \stepcounter{utterance}  
    & & & \multicolumn{2}{p{0.3\linewidth}}{
        \cellcolor[rgb]{0.9,0.9,0.9}{
            \makecell[{{p{\linewidth}}}]{
                \texttt{\tiny{[GM$|$GM]}}
                \texttt{game\_result = LOSE} \\
            }
        }
    }
    & & \\ \\

    \theutterance \stepcounter{utterance}  
    & & \multicolumn{4}{p{0.6\linewidth}}{
        \cellcolor[rgb]{0.95,0.95,0.95}{
            \makecell[{{p{\linewidth}}}]{
                \texttt{\tiny{[P1$\leftrightarrow$P]}}
                \texttt{Post{-}episode Reflexion (not sent to GM):} \\
\texttt{New plan: In the previous attempt, I correctly identified the error as a change in the adjective "unvollständige" (incomplete) to "complete" and selected Option (A) "Modifiers or Adjectives". However, the outcome was "LOSE", indicating that my classification was incorrect according to the specific logic of this task. Re{-}evaluating the definitions: "Negation or Antonyms" is defined as changing comparatives to their antonyms. The shift from "incomplete" to "complete" is a direct antonym swap of the adjective describing the list. In the third example, "is not" vs "is" was classified as Negation/Antonyms. Similarly, flipping an adjective to its direct opposite (incomplete {-}> complete) fundamentally alters the truth value of the statement in a way that aligns with "Negation or Antonyms" rather than just a generic modifier change. My mistake was over{-}prioritizing the "Modifiers" category for adjective changes, whereas the specific nature of the change (an antonym swap) places it in the "Negation or Antonyms" category. Next time, when an adjective is changed to its direct opposite (e.g., incomplete to complete, small to large), I will classify it as (C) Negation or Antonyms, as this represents a logical negation of the original attribute.} \\
            }
        }
    }
    & \\ \\

\end{supertabular}
}

\end{document}
